\section[Cote 90 : un groupe plat irréductible dont un saturé n'est pas de type fini]{Cote 90 : un groupe plat irréductible dont un saturé n'est pas de type fini\protect\verifie{Folder90}}\label{sec:90}

La cote 90, \emph{KG (X)}, est datée par l'inventaire \q{à partir de 1964} ; ses notes sont écrites au dos de feuillets dactylographiés d'EGA~IV, §\,18. Elles préparent l'existence de résolutions équivariantes. La page~1 pose, sur un anneau de Dedekind $A$, la notion de \emph{saturé} : pour un sous-module $N$ d'un module sans torsion $M$, le saturé $\widetilde N$ est l'ensemble des $y \in M$ dont un multiple non nul est dans $N$. Le lemme~1 donne des conditions équivalentes pour que le saturé de tout sous-module de type fini soit de type fini (\lot{90}{1}{1}). La page~9 applique ces conditions à l'algèbre $B$ d'un groupe affine et plat $G$ sur $A$, et dit dans une parenthèse (\lot{90}{1}{9}) :

\begin{quote}
\emph{cela signifie que pour les composantes irréd.\ de G, leurs fibres sont non vides}
\end{quote}

\noindent Pour un groupe de type fini, la parenthèse est raisonnable. Sans cette hypothèse, que la proposition ne pose pas, elle tombe en défaut.

\noindent\emph{Le contre-exemple.} Soient $R$ un anneau de valuation discrète, $K$ son corps des fractions, $\varpi$ une uniformisante et $k = R/\varpi R$ le corps résiduel. On pose
\[
B = \{ f \in K[T] \;:\; f(0) \in R \},
\]
les polynômes à coefficients dans $K$ dont le terme constant est entier. C'est la réunion des dilatations successives $R[T/\varpi^n]$ du groupe additif en l'origine : on autorise des dénominateurs arbitraires partout, sauf au point $T = 0$ de la fibre spéciale. Cet exemple figure déjà dans SGA~3 (exposé~VI\textsubscript{B}, 13.3), où l'anneau est décrit comme $R \oplus K \oplus K \oplus \cdots$, non projectif ; Duong, Hai et dos Santos le rappellent (arXiv:1701.06518, exemple~1.1). La page de la cote~90, qui ne suppose pas $G$ de type fini, le contredit donc dans le séminaire même de Grothendieck.

\begin{encadre}
\begin{contrexemple}[SGA~3, exposé~VI\textsubscript{B}, 13.3, contre la parenthèse de la page 9 \cite{sga3-1}]
L'anneau $B$ est intègre, de sorte que $\mathrm{Spec}\,B$ est irréductible ; il est plat sur $R$ ; l'application $f \mapsto f(0) \bmod \varpi$ est surjective de $B$ sur $k$, de noyau $\varpi B$, de sorte que $B/\varpi B \simeq k$ ; et $K \otimes_R B \neq 0$. Les deux fibres sont donc non vides. Pour tout $f \in B$, il existe $b_i, c_i \in B$ en nombre fini tels que $f(x + y) = \sum_i b_i(x)\,c_i(y)$ pour tous $x, y \in K$ : la comultiplication du groupe additif préserve $B$. Pourtant le saturé de $R \cdot T$ dans $B$ est $K \cdot T$, qui n'est pas un $R$-module de type fini.
\end{contrexemple}
\end{encadre}

\begin{proof}
$B$ est un sous-anneau de $K[T]$, donc intègre. Il est sans torsion : si $r \neq 0$ et $r f = 0$ dans $K[T]$, alors $f = 0$, puisque $r$ est inversible dans $K$. Sur un anneau principal, un module sans torsion est plat.

Soit $f \in B$ avec $f(0) \in \varpi R$, disons $f(0) = a\varpi$. Alors $g = \varpi^{-1} f$ a pour terme constant $a \in R$, donc $g \in B$ et $f = \varpi g \in \varpi B$. Réciproquement, un élément de $\varpi B$ a un terme constant dans $\varpi R$. Le noyau de $f \mapsto f(0) \bmod \varpi$ est donc $\varpi B$, et l'application est surjective puisque les constantes de $R$ sont dans $B$. Pour la fibre générique, l'application $K \otimes_R B \to K[T]$, $\lambda \otimes f \mapsto \lambda f$, envoie $1 \otimes 1$ sur $1 \neq 0$.

Pour la comultiplication, écrivons $f = \sum_{k \leq d} f_k T^k$ avec $f_0 \in R$. La formule du binôme donne $f(x + y) = \sum_k \sum_{j \leq k} \binom{k}{j} f_k\, x^j y^{k-j}$. On répartit chaque terme en un produit de deux éléments de $B$ : pour $k = j = 0$, $f_0 \cdot 1$ ; pour $j = 0 < k$, le facteur $1$ en $x$ et $f_k y^k$ en $y$, de terme constant nul ; pour $j \geq 1$, le facteur $\binom{k}{j} f_k x^j$, de terme constant nul, et $y^{k-j}$. Tous ces facteurs sont dans $B$.

Le saturé. Si $r \neq 0$ et $r f = a T$ avec $a \in R$, alors $f = (a/r) T \in K \cdot T$. Inversement, pour $c \in K$, écrivons $c = a/s$ avec $a, s \in R$, $s \neq 0$ ; alors $s \cdot cT = aT \in R \cdot T$. Le saturé est donc $K \cdot T$, contenu dans $B$ puisque $cT$ s'annule en~$0$. L'application $cT \mapsto c$ est un isomorphisme de $R$-modules de $K \cdot T$ sur $K$. Si $K$ était de type fini sur $R$, il serait entier sur $R$ ; or $R$ est intégralement clos, donc $\varpi^{-1}$ serait dans $R$, et $\varpi$ serait inversible, ce qui n'est pas.
\end{proof}

\begin{remarque}
Le lemme~1 de la page~1 dit que, pour un module sans torsion sur un anneau de Dedekind, les saturés des sous-modules de type fini sont de type fini si et seulement si le module est réunion filtrante de sous-modules de type fini saturés. Ici $R \cdot T$ est de type fini et son saturé ne l'est pas : $B$ ne satisfait pas les conditions du lemme~1, alors que $\mathrm{Spec}\,B$, qui n'a qu'une composante irréductible, a ses deux fibres non vides. Que $B$ soit une algèbre de Hopf n'est vérifié que par l'identité de comultiplication, point par point sur $K$ ; comme $K$ est infini, c'est une identité de polynômes. L'antipode, $f(T) \mapsto f(-T)$, et la counité, $f \mapsto f(0)$, préservent visiblement $B$, mais ne sont pas écrites en Lean.
\end{remarque}

\noindent\emph{Le cas de type fini et réduit.} Pour un groupe de type fini, la parenthèse se ramène à une question d'algèbre commutative sur l'idéal $d(A) = \bigcap_n \pi^n A$ des éléments infiniment divisibles par l'uniformisante. On en démontre ici le cœur, valable pour tout élément $\pi$ d'un anneau commutatif.

\begin{encadre}
\begin{enonce}[Composantes et éléments infiniment divisibles, énoncé de l'édition]\verifie[\leanblobmain]{Folder90}
Soient $A$ un anneau commutatif, $\pi \in A$ et $d(A) = \bigcap_{n} \pi^n A$.
\begin{enumerate}
\item[(a)] Si un idéal premier minimal $\mathfrak p$ de $A$ vérifie $\mathfrak p + \pi A = A$, autrement dit si la composante irréductible $V(\mathfrak p)$ de $\mathrm{Spec}\,A$ ne rencontre pas la fibre spéciale $V(\pi)$, alors $d(A) \neq 0$.
\item[(b)] Si $A$ est noethérien et réduit et si tout premier minimal $\mathfrak p$ vérifie $\mathfrak p + \pi A \neq A$, alors $d(A) = 0$.
\end{enumerate}
Pour $A$ noethérien et réduit, $d(A) = 0$ si et seulement si toutes les composantes irréductibles de $\mathrm{Spec}\,A$ rencontrent la fibre spéciale.
\end{enonce}
\end{encadre}

\begin{proof}
(a) Comme $\mathfrak p + \pi A = A$, il existe $c \in A$ tel que $1 - c\pi \in \mathfrak p$. Dans un anneau commutatif, tout élément d'un premier minimal est diviseur de zéro : il existe $a \neq 0$ tel que $a(1 - c\pi) = 0$. Alors $a = ac\pi$ et, par récurrence, $a = a(c\pi)^n = (ac^n)\,\pi^n$ pour tout $n$ : $a$ est un élément non nul de $d(A)$. L'idée : sur une composante qui évite la fibre spéciale, $\pi$ est inversible, et une fonction portée par cette composante est divisible par toutes les puissances de $\pi$. Cette partie n'utilise ni l'hypothèse noethérienne ni le théorème de Krull.

(b) Soit $\mathfrak p$ un premier minimal. L'anneau $A/\mathfrak p$ est intègre et noethérien, et l'image de $\pi$ y engendre un idéal propre, puisque $\mathfrak p + \pi A \neq A$. Le théorème d'intersection de Krull donne $\bigcap_n \bar\pi^n (A/\mathfrak p) = 0$ ; l'image de $d(A)$ dans $A/\mathfrak p$ est donc nulle, et $d(A) \subseteq \mathfrak p$. L'intersection des premiers minimaux est le nilradical, nul puisque $A$ est réduit : $d(A) = 0$.
\end{proof}

\noindent Lean : \href{\leanblobmain/Folder90.lean}{\texttt{iInf\_pow\_ne\_bot\_of\_minimalPrime}} pour (a), \href{\leanblobmain/Folder90.lean}{\texttt{iInf\_pow\_eq\_bot\_of\_minimalPrimes}} pour (b) et \href{\leanblobmain/Folder90.lean}{\texttt{iInf\_pow\_eq\_bot\_iff}} pour l'équivalence.

\begin{remarque}
Pour $A = R[G]$, algèbre d'un groupe affine et plat de type fini sur un anneau de valuation discrète complet $R$ d'uniformisante $\pi$, la proposition~6.19 de Hai et dos Santos \cite{haidossantos2018} et le théorème de Kaplansky sur les modules projectifs traduisent $d(A) = 0$ en deux propriétés : $G$ est \emph{spécialement localement fini}, et $R[G]$ est un $R$-module libre. Avec l'énoncé ci-dessus, on obtient le critère de la page~9 quand $R[G]$ est réduit, ce qui est toujours le cas en caractéristique générique nulle. Ces deux résultats sont cités, non formalisés. Le cas non réduit, en caractéristique générique $p$, reste ouvert.
\end{remarque}

Ne sont pas démontrés ici : le lemme~1 lui-même, et l'énoncé suivant de la cote.

\begin{encadre}
\begin{theoreme*}[D'après les pages 13 à 23, sous-modules équivariants sur une base quelconque, forme précisée par la lecture\piste{90-equivariant-resolution-general-base}]\nonverifie
Soient $S$ un schéma et $G$ un $S$-schéma en groupes plat et affine, dont l'algèbre est réunion filtrante de sous-cogèbres localement libres de rang fini à quotients plats, opérant sur un $S$-schéma $X$ quasi-compact et quasi-séparé muni d'un faisceau ample relatif $G$-linéarisé. Tout sous-module quasi-cohérent d'un $G$-module quasi-cohérent $E$ sur $X$ est contenu dans un plus petit sous-module quasi-cohérent $G$-stable, compatible au changement de base plat sur $S$ et de type fini quand le premier l'est.
\end{theoreme*}
\end{encadre}
{\small Les feuillets en esquissent une démonstration incomplète (pages 19 à 23). L'argument s'arrête sur \q{Rédiger mieux ce début}. Mathlib n'a pas les modules quasi-cohérents équivariants. On se reportera aux notes de Grothendieck (\lot{90}{1}{13}).\par}

\subsection*{Le plus petit sous-comodule, sur $\mathrm{Spec}\,A$ (pages 7--9 et 13--15)\protect\verifie{Folder90Comodules}}\label{sec:90-comodules}

La page~13 (\lot{90}{1}{13}) énonce, pour un groupe affine sur une base quelconque dont l'algèbre $B$ est réunion filtrante de sous-modules $B_i$ avec 1)~$B/B_i$ plat, 2)~$\Delta B_i \subset B_i \otimes B_i$, 3)~$B_i$ localement libre de type fini, et pour un $G$-module quasi-cohérent $E$ à action triviale :

\begin{quote}
\emph{1) tout sous-module quasi-cohérent $F$ de $E$ est contenu dans un plus petit sous-module quasi-cohérent $F'$ stable par $G$ ; sa formation commute à tout changement de base plat ; si $F$ est de type fini, $F'$ aussi ; 2) si $X$ est quasi-compact et quasi-séparé, $E$ est réunion filtrante de sous-modules quasi-cohérents de type fini stables par $G$.}
\end{quote}

\noindent La page n'en donne pas de démonstration ; la lecture propose la formule $F' = \sum_i U_i \cdot (F \cap E_i)$, $U_i = \check B_i$, tirée de la description des pages~7 et~9. On démontre ici le cas de l'anneau, $X = S = \mathrm{Spec}\,A$, que la page énonce et argumente. Mathlib a les cogèbres mais pas les comodules : un comodule est un $A$-module $E$ muni d'une coaction $\delta \colon E \to B \otimes E$ counitaire et coassociative, et $N \subset E$ est un sous-comodule si $\delta N$ est contenu dans l'image de $B \otimes N$.

\begin{encadre}
\begin{enonce}[Pages 13--15, cas affine]\verifie[\leanblobmain]{Folder90Comodules}
Soient $A$ un anneau commutatif et $B$ une $A$-cogèbre, réunion filtrante de sous-modules $B_i$ projectifs de type fini, avec $B/B_i$ plat et $\Delta B_i \subset \mathrm{Im}(B_i \otimes B_i \to B \otimes B)$. Soit $E$ un $B$-comodule.
\begin{enumerate}
  \item Tout sous-module $F$ de $E$ est contenu dans un plus petit sous-comodule $F'$, à savoir le sous-module engendré par les $(u \otimes \mathrm{id})(y)$, où $y \in B_i \otimes E$ relève $\delta f$, $f \in F$ et $u \in U_i = \mathrm{Hom}_A(B_i, A)$.
  \item Si $F$ est de type fini, $F'$ l'est, et $F' = U_i \cdot F$ pour un seul indice $i$.
  \item Tout élément de $E$ est dans un sous-comodule de type fini ; ceux-ci sont stables par somme, et $E$ en est la réunion filtrante.
  \item Pour tout morphisme de comodules $\varphi$, $\varphi(F') = (\varphi F)'$.
  \item Pour toute $A$-algèbre $A'$, $A' \otimes E$ est un comodule sur la $A'$-cogèbre $A' \otimes B$, et l'image de $A' \otimes F'$ dans $A' \otimes E$ est le plus petit sous-comodule contenant l'image de $A' \otimes F$. Si $A'$ est plat sur $A$, cette image est $A' \otimes F'$.
\end{enumerate}
\end{enonce}
\end{encadre}

\begin{proof}
\emph{L'étape de platitude.} Comme $B/B_i$ est plat, $B_i \otimes M \to B \otimes M$ est injectif pour tout $M$ (suite exacte $0 \to B_i \to B \to B/B_i \to 0$, lemme du serpent sur une présentation de $M$). Il en découle l'égalité de la lecture, $(B_i \otimes E) \cap (B \otimes N) = B_i \otimes N$ : si $y \in B_i \otimes E$ a son image dans $B \otimes N$, son image dans $B_i \otimes (E/N)$ s'envoie sur $0$ dans $B \otimes (E/N)$, donc est nulle, et $y$ provient de $B_i \otimes N$ par exactitude à droite.

\emph{$F'$ contient $F$.} Pour $f \in F$, $\delta f$ est une somme finie, donc dans un $B_i \otimes E$ ; la coünité restreinte à $B_i$ donne $(\varepsilon \otimes \mathrm{id})(y) = (\varepsilon \otimes \mathrm{id})\,\delta f = f$.

\emph{Minimalité.} Si $N \supset F$ est un sous-comodule et $y \in B_i \otimes E$ relève $\delta f$, $f \in F$, l'image de $y$ est dans $B \otimes N$, donc $y \in B_i \otimes N$ par l'étape de platitude, et $(u \otimes \mathrm{id})(y) \in N$.

\emph{Stabilité.} Le module projectif de type fini $B_i$ a une base duale finie $(b_k, b_k^*)$, et $\Delta$ se relève en $\Delta_i \colon B_i \to B_i \otimes B_i$ par projectivité. La coassociativité et l'injectivité de $B_i \otimes (B \otimes E) \to B \otimes (B \otimes E)$ donnent $(\mathrm{id} \otimes \delta)(y) = (\Delta_i \otimes \mathrm{id})(y)$ dans $B_i \otimes B \otimes E$, d'où
\[
\delta\big((u \otimes \mathrm{id})\,y\big) = \sum_k b_k \otimes \big((u * b_k^*) \otimes \mathrm{id}\big)(y), \qquad u * v = (u \otimes v) \circ \Delta_i ,
\]
et les coefficients sont encore des générateurs de $F'$.

\emph{Finitude.} Si $F$ est engendré par $f_1, \dots, f_m$, on choisit un $i$ commun (la famille est filtrante) et des relèvements $y_j$ de $\delta f_j$. Le sous-module $G$ engendré par les $(b_k^* \otimes \mathrm{id})(y_j)$ contient tous les $(u \otimes \mathrm{id})(y_j)$, puisque $u = \sum_k u(b_k)\, b_k^*$ ; il est stable par la formule précédente et contient $F$ ; donc $F' \subset G \subset F'$, et $F'$ est de type fini. Le point~3 s'en déduit avec $F = A e$.

\emph{Morphismes et changement de base.} Un morphisme $\varphi$ envoie les générateurs de $F'$ sur des générateurs de $(\varphi F)'$, et $\varphi(F')$ est un sous-comodule contenant $\varphi F$. Pour $A'$, l'application $e \mapsto 1 \otimes e$ est un morphisme de $B$-comodules $E \to A' \otimes E$ ; un sous-comodule de $A' \otimes E$ sur $A' \otimes B$ en est un sur $B$, par $(a \otimes b) \otimes n \mapsto b \otimes a n$ ; et l'image de $A' \otimes F'$ est le $A'$-module engendré par $1 \otimes F'$. Si $A'$ est plat, $A' \otimes F' \to A' \otimes E$ est injectif.
\end{proof}

\begin{remarque}
L'énoncé~1 tient sur tout anneau commutatif tel que la lecture le complète, et la compatibilité au changement de base tient pour toute $A$-algèbre $A'$ dès qu'on l'énonce avec des images ; la platitude de $A'$ ne sert qu'à identifier l'image à $A' \otimes F'$. La platitude de $B/B_i$ n'intervient que par l'injectivité de $B_i \otimes M \to B \otimes M$. Ne sont pas formalisés : les modules quasi-cohérents sur un schéma $X$ au-dessus d'une base $S$ quelconque (recollement, hypothèse « quasi-compact et quasi-séparé » de 2)), les énoncés~3) et~4), l'équivalence de la page~7 entre comodules $i$-admissibles et $U_i$-modules, et l'action de $G$ sur $X$ des pages~17 à~23, qui est l'objet de la trouvaille. Noms Lean : \texttt{Comodule}, \texttt{IsSubcomodule}, \texttt{Exhaustion}, \texttt{Exhaustion.closure}, \texttt{Exhaustion.isLeast\_closure}, \texttt{Exhaustion.closure\_fg}, \texttt{Exhaustion.exists\_fg\_isSubcomodule}, \texttt{IsSubcomodule.sup}, \texttt{Exhaustion.map\_closure}, \texttt{Exhaustion.coact\_coeff}, \texttt{Exhaustion.mem\_range\_lTensor\_of\_rTensor\_mem}, \texttt{comoduleBaseChange}, \texttt{isLeast\_closure\_baseChange}, \texttt{baseChange\_closure\_injective}.
\end{remarque}
